% Generated draft: requires mathematical and source review.
\documentclass[11pt]{article}
\usepackage[margin=25mm]{geometry}
\usepackage{fontspec}
\setmainfont{Latin Modern Roman}
\usepackage{amsmath,amssymb,mathtools,unicode-math}
\setmathfont{Latin Modern Math}
\usepackage{enumitem,needspace,xurl,hyperref}
\setlength{\parindent}{0pt}
\setlength{\parskip}{5pt}
\setlength{\emergencystretch}{3em}
\newcommand{\N}{\mathbb{N}}
\newcommand{\Z}{\mathbb{Z}}
\newcommand{\Q}{\mathbb{Q}}
\newcommand{\R}{\mathbb{R}}
\newcommand{\C}{\mathbb{C}}
\newcommand{\F}{\mathbb{F}}
\newcommand{\A}{\mathbb{A}}
\newcommand{\PP}{\mathbb{P}}
\newcommand{\E}{\mathbb{E}}
\newcommand{\eps}{\varepsilon}
\newcommand{\dd}{\,\mathrm{d}}
\begin{document}
\section*{UnsolvedMath ID: 30005470}
\textbf{Source problem code:} catalogue entry
\subsection{Definitions and mathematical statement}
Let \(n\geq 1\) be an integer. A function \(f:\R^n\to\R\) is real analytic if it is locally represented by convergent real power series, and it is positive semidefinite if \(f(x)\geq 0\) for every \(x\in\R^n\).

Under the local-quotient convention, a real-analytic meromorphic function on \(\R^n\) is a global meromorphic section locally represented by quotients \(a/b\), where \(a,b\) are real-analytic functions and the denominator has nonzero germ. Local representatives agree wherever their quotients are defined. An identity between such functions is understood as a meromorphic identity, checked locally after clearing denominators; poles of individual summands are permitted.

The supplied question, with ``sum of squares'' interpreted as a finite sum, is whether
\[
\forall n\geq 1\;\forall f\in C^\omega(\R^n,\R),\qquad
\bigl(\forall x\in\R^n,\ f(x)\geq 0\bigr)
\Longrightarrow
\bigl(\exists m\geq 1\;\exists g_1,\ldots,g_m\text{ real-analytic meromorphic on }\R^n,\quad f=\sum_{j=1}^m g_j^2\bigr).
\]
Here \(m\) may depend on \(n\) and \(f\). No compactness assumption or bound on the number of squares is supplied.

The source does not define ``real-analytic meromorphic.'' If it requires each \(g_j\) to be a quotient of globally defined real-analytic functions, the corresponding precise formulation is: for every such \(n\) and \(f\), do there exist \(m\geq 1\) and \(h,a_1,\ldots,a_m\in C^\omega(\R^n,\R)\), with \(h\not\equiv 0\), such that
\[
h^2f=\sum_{j=1}^m a_j^2\quad\text{on }\R^n?
\]
This common-denominator formulation expresses the global-quotient convention. The supplied evidence does not establish its equivalence with the local-quotient formulation, so the intended convention requires review.
\subsection{Short English statement}
Does every everywhere nonnegative real-analytic function on \(\R^n\), for every positive integer \(n\), admit a finite representation as a sum of squares of real-analytic meromorphic functions? The identity allows poles and is interpreted after clearing denominators. The source leaves the precise meromorphic convention unstated.
\subsection{Sources}
\begin{itemize}
\item Supplied local UnsolvedMath problem registry, permanent ID \(30005470\), exact source code \texttt{catalogue entry}, title ``Analytic Hilbert Seventeenth Problem,'' in the Oberwolfach Reports Open Problems collection. No individual proposer is identified; the supplied proposed year is 2023.
\item \emph{New Directions in Real Algebraic Geometry} (2023), Oberwolfach Reports, pp.\ 813--881, DOI \texttt{10.4171/owr/2023/15}: \url{https://doi.org/10.4171/owr/2023/15}. The supplied passage labels the question ``Open question 1'' and describes the real-analytic analogue of Artin's theorem as open. The registry labels the statement \texttt{corrected\_verified} and reports that it was reconstructed and checked against the official report; this is supplied verification metadata, not an independent check performed for this draft.
\item Supplied literature reference: Fernando, ``meromorphic functions on Stein surfaces,'' \url{https://doi.org/10.1515/crelle-2025-0071}.
\item Supplied literature reference: Fernando, ``positive semidefinite analytic functions on surfaces,'' \url{https://doi.org/10.1007/s12220-021-00721-1}.
\item Literature assessment dated exactly \texttt{2026-08-22}: the supplied curation labels the problem partially solved and states that the analytic Hilbert-seventeenth analogue is known in dimensions at most two and under compactness hypotheses, that recent work removes compactness for Stein surfaces, and that the global higher-dimensional statement remains unresolved. This is supplied dated curation derived from the OWR material, not independently verified here; its time-sensitive status and cited claims require independent verification. The supplied sources checked include both Fernando references above and the OWR report at \url{https://doi.org/10.4171/owr/2023/15}.
\end{itemize}
\end{document}
